% !TEX root = ../main.tex
\begin{ouverture}[ensembles]{13}{Ensembles et raisonnement}{Partie 5 · Méthodes transversales}
{Ensembles\enspace·\enspace Logique\enspace·\enspace Quantificateurs\enspace·\enspace Démonstrations}
\objectif{Utiliser les opérations sur les ensembles et les intervalles.}
\objectif{Distinguer implication, réciproque et contraposée.}
\objectif{Nier une proposition et identifier les rôles des lettres.}
\objectif{Rédiger des preuves par cas, par contraposée ou par l’absurde.}
\prerequis{Ensembles de nombres, calcul littéral et intervalles.}
\prerequis{Parité des entiers et fonctions du second degré.}
\end{ouverture}
\essentiel
\essSection{1}{Ensembles, intervalles et couples}
\begin{definition}
$x\in A$ signifie que $x$ est un élément de $A$ ; $A\subseteq B$ que tout
élément de $A$ appartient à $B$. L'ensemble vide est noté $\varnothing$.
$A\cap B$ contient les éléments communs ; $A\cup B$ ceux de l'un ou de l'autre.
Dans un univers $E$, le complémentaire est $E\setminus A$, aussi noté $\overline A$.
Le cardinal $\operatorname{Card}(A)$ compte les éléments d'un ensemble fini.
\end{definition}
\begin{propriete}
Un couple est ordonné : $(a;b)$ et $(b;a)$ ne sont pas en général identiques.
Le produit cartésien $A\times B$ est l'ensemble des couples $(a;b)$ avec
$a\in A$ et $b\in B$. Si les deux ensembles sont finis,
$\operatorname{Card}(A\times B)=\operatorname{Card}(A)\operatorname{Card}(B)$.
\end{propriete}
\begin{center}
\begin{tikzpicture}[scale=.8]
\draw[otMuted,rounded corners](0,0)rectangle(7,3);
\fill[otNavy,opacity=.12](2.5,1.5)circle(1.2);
\fill[otGreen,opacity=.2](4.2,1.5)circle(1.2);
\draw[otNavy](2.5,1.5)circle(1.2);\draw[otGreenDark](4.2,1.5)circle(1.2);
\node at(2,1.5){$A$};\node at(4.8,1.5){$B$};\node[above]at(3.35,1.5){$A\cap B$};
\node at(.4,2.6){$E$};
\end{tikzpicture}
\par\legende{La réunion inclut les deux disques ; l’intersection est leur zone commune.}
\end{center}

\begin{attention}
Ne pas confondre appartenance et inclusion : $2\in\{2;4\}$ mais
$\{2\}\subseteq\{2;4\}$. Une borne exclue d'un intervalle n'en est pas un élément.
\end{attention}

\clearpage
\essSection{2}{Implication, réciproque et conditions}
\begin{definition}
L'implication « si $P$, alors $Q$ » affirme que $Q$ est vraie dès que $P$
est vraie. Sa réciproque est « si $Q$, alors $P$ ».
Sa contraposée est « si non $Q$, alors non $P$ » ; elle lui est équivalente.
Une équivalence réunit une implication et sa réciproque.
\end{definition}
\begin{propriete}
Dans $P\Rightarrow Q$, $P$ est une condition suffisante pour $Q$ ; $Q$ est
une condition nécessaire pour $P$. Exemple : être multiple de 4 suffit
pour être pair ; être pair est nécessaire pour être multiple de 4.
\end{propriete}
\essSection{3}{Quantifier, nier et démontrer}
\begin{propriete}
La négation de « pour tout $x$, $P(x)$ » est « il existe $x$ pour lequel $P(x)$
est fausse ». La négation de « il existe $x$ tel que $P(x)$ » est
« pour tout $x$, $P(x)$ est fausse ». Nier « $P$ et $Q$ » donne « non $P$ ou non $Q$ ».
Nier « $P$ ou $Q$ » donne « non $P$ et non $Q$ ».
\end{propriete}
\begin{savoirfaire}
Une preuve directe enchaîne des propriétés connues. Une preuve par cas
traite toutes les possibilités. Une preuve par contraposée démontre la
contraposée. Une preuve par l'absurde suppose le contraire de la conclusion
et obtient une contradiction. Un contre-exemple suffit à réfuter une affirmation universelle.
\end{savoirfaire}
\begin{attention}
Une identité est vraie pour toutes les valeurs autorisées ; une équation
demande de chercher des valeurs de l'inconnue. Dans une famille $f_m(x)$,
$x$ peut être une variable et $m$ un paramètre. Un test sur quelques nombres
ne démontre pas une affirmation « pour tout ».
\end{attention}

\exercices

\rubrique{Découverte}
\begin{blocLivre}{ex-13-1}
\begin{exercice}[1]{Appartenance et opérations}
On donne $E=\{1;2;3;4;5;6\}$, $A=\{2;4;6\}$ et $B=\{1;2;3\}$.\begin{questions}
\item Déterminer $A\cap B$, $A\cup B$ et $E\setminus A$.
\item Corriger si nécessaire : $2\subseteq A$ et $\{2\}\in A$.
\end{questions}
\end{exercice}
\end{blocLivre}
\begin{blocLivre}{ex-13-2}
\begin{exercice}[1]{Réunir des intervalles}
On donne $A=[-2;3]$ et $B=]1;5[$.\begin{questions}
\item Déterminer $A\cap B$ et $A\cup B$.
\item Donner le complémentaire de $A$ dans $\R$.
\end{questions}
\label{figlivre-r13-02-intervalles-ex}
\figureLivre{r13-02-intervalles}

\end{exercice}
\end{blocLivre}
\begin{blocLivre}{ex-13-3}
\begin{exercice}[1]{Construire des couples}
On pose $A=\{a;b\}$ et $B=\{0;1;2\}$.\begin{questions}
\item Lister $A\times B$ et donner son cardinal.
\item Le couple $(1;a)$ appartient-il à $A\times B$ ? À $B\times A$ ?
\end{questions}
\end{exercice}
\end{blocLivre}
\begin{blocLivre}{ex-13-4}
\begin{exercice}[1]{Trouver un contre-exemple}
Réfuter les deux affirmations suivantes :\begin{questions}
\item « Pour tout réel $x$, $x^2\geq x$. »
\item « Si $x^2=4$, alors $x=2$. »
\end{questions}
\end{exercice}
\end{blocLivre}
\begin{blocLivre}{ex-13-5}
\begin{exercice}[1]{Nier une phrase quantifiée}
Pour chacune des phrases suivantes, écrire sa négation puis donner la valeur de vérité des deux phrases :\begin{questions}
\item « Pour tout réel $x$, $x^2\geq0$. »
\item « Il existe un réel $x$ tel que $x^2=-1$. »
\end{questions}
\end{exercice}
\end{blocLivre}
\clearpage
\rubrique{Entraînement}
\begin{blocLivre}{ex-13-6}
\begin{exercice}[2]{Réciproque et contraposée}
On considère l’implication « si $x>2$, alors $x^2>4$ », pour $x$ réel.\begin{questions}
\item Justifier cette implication.
\item Écrire sa réciproque et décider si elle est vraie.
\item Écrire sa contraposée et sa valeur de vérité.
\end{questions}
\end{exercice}
\end{blocLivre}
\begin{blocLivre}{ex-13-7}
\begin{exercice}[2]{Nécessaire ou suffisante}
Pour un entier $n$, comparer les propriétés P : « $n$ est multiple de 4 » et Q : « $n$ est pair ».\begin{questions}
\item Quelle implication est toujours vraie ?
\item Compléter : P est une condition \ldots\ pour Q ; Q est une condition \ldots\ pour P.
\item Donner un exemple montrant que l’équivalence est fausse.
\end{questions}
\end{exercice}
\end{blocLivre}
\begin{blocLivre}{ex-13-8}
\begin{exercice}[2]{Variable, inconnue et paramètre}
\begin{questions}
\item Quel est le statut de $x$ dans l’identité $x^2-1=(x-1)(x+1)$ ?
\item Dans l’équation $x^2-1=0$, quel est le rôle de $x$ et quelles sont ses valeurs solutions ?
\item Dans $f_m(x)=x^2-mx+1$, distinguer variable et paramètre.
\end{questions}
\end{exercice}
\end{blocLivre}
\begin{blocLivre}{ex-13-9}
\begin{exercice}[2]{Complémentaires et connecteurs}
Dans $E=\{1;\ldots;8\}$, A est l’ensemble des entiers pairs et B celui des entiers inférieurs ou égaux à 4.\begin{questions}
\item Calculer $E\setminus(A\cup B)$.
\item Calculer $(E\setminus A)\cap(E\setminus B)$ et comparer.
\item Expliquer l’égalité obtenue en utilisant « et » et « ou ».
\end{questions}
\end{exercice}
\end{blocLivre}
\begin{blocLivre}{ex-13-10}
\begin{exercice}[2]{Lire une quantification implicite}
On rappelle : « f est croissante sur I » signifie que pour tous $a,b\in I$ tels que $a<b$, $f(a)\leq f(b)$.\begin{questions}
\item Écrire la négation de cette définition.
\item Montrer que $f(x)=x^2$ n’est pas croissante sur $\R$.
\item Montrer directement qu’elle est croissante sur $[0;+\infty[$.
\end{questions}
\end{exercice}
\end{blocLivre}
\begin{blocLivre}{ex-13-11}
\begin{exercice}[2]{Démontrer par disjonction des cas}
\begin{questions}
\item Montrer que $n(n+1)$ est pair pour tout entier naturel $n$, en distinguant $n$ pair et $n$ impair.
\item Montrer que $|x|\geq x$ pour tout réel $x$, en distinguant le signe de $x$.
\end{questions}
\end{exercice}
\end{blocLivre}
\clearpage
\rubrique{Problèmes}
\begin{blocLivre}{ex-13-12}
\begin{exercice}[3]{Une preuve par contraposée}
On veut démontrer : « si le carré d’un entier est impair, alors cet entier est impair ».\begin{questions}
\item Écrire la contraposée.
\item La démontrer en écrivant un entier pair sous la forme $2k$.
\item Démontrer aussi la réciproque et formuler l’équivalence obtenue.
\end{questions}
\end{exercice}
\end{blocLivre}
\begin{blocLivre}{ex-13-13}
\begin{exercice}[3]{Une preuve par l'absurde}
On veut montrer que $\sqrt2$ n’est pas rationnel. On suppose au contraire $\sqrt2=p/q$, avec $p,q$ entiers positifs sans facteur commun.\begin{questions}
\item En déduire $p^2=2q^2$ puis montrer que $p$ est pair.
\item Écrire $p=2k$ et montrer que $q$ est pair.
\item Identifier la contradiction et conclure.
\end{questions}
\end{exercice}
\end{blocLivre}
\begin{blocLivre}{ex-13-14}
\begin{exercice}[3]{Une transformation sous condition}
Résoudre $\dfrac{x^2-1}{x-1}=2$.\begin{questions}
\item Préciser le domaine et simplifier sous cette condition.
\item Résoudre l’équation simplifiée puis conclure.
\item Expliquer l’erreur d’une résolution qui multiplierait sans conserver la valeur interdite.
\end{questions}
\end{exercice}
\end{blocLivre}
\begin{blocLivre}{ex-13-15}
\begin{exercice}[3]{Quantifier une famille de trinômes}
Pour un réel $m$, on pose $P_m(x)=x^2-2mx+1$.\begin{questions}
\item Écrire $P_m$ sous forme canonique.
\item Déterminer les $m$ tels que « pour tout réel $x$, $P_m(x)\geq0$ ».
\item Déterminer les $m$ tels que « pour tout réel $x$, $P_m(x)>0$ ».
\item Écrire la négation de la phrase de la question 2 et trouver les $m$ correspondants.
\item Présenter oralement, en deux minutes, la distinction entre les questions 2 et 3.
\end{questions}
\end{exercice}
\end{blocLivre}
